\chapter{Conclusion}

In this work, the \gls{x-bap} method is presented, which is a \gls{psf}-corrected aperture photometry. The goal of \gls{x-bap} is to accurately measure colors across different photometric bands, compensating for the effects that the \gls{psf} and seeing have on the observation. The improved color measurements can be used for applications such as source classification and measuring photometric redshift. The main findings in this work are as follows:
\begin{itemize}
    \item \gls{x-bap} only works when the pre-seeing aperture is bigger than the largest \gls{psf} when applied to multiple bands.
    \item The error on the aperture flux measurement of \gls{x-bap} can be accurately computed with sufficient information about the noise in the observation.
    \item Using different aperture sizes with \gls{x-bap} can give information about the shape of the source. The flux measured from point sources does not change with the size of the Gaussian aperture when it is centered on the source, whereas it does for extended sources.
    \item Elliptical Gaussian apertures only provide limited improvement of the \gls{snr} compared to circular Gaussian apertures.
    \item Applying \gls{x-bap} to observations by Euclid shows similar colors compared to the colors derived from the flux in the \gls{mer} catalogs, which shows that \gls{x-bap} does perform correctly
    \item \gls{x-bap} can be applied to observations for classification, but deblending is an important part to accurately classify sources.
\end{itemize}
As more and more large sky surveys by telescopes such as Euclid and Vera Rubin Observatory produce more multi-band data than ever, it becomes increasingly important to combine these accurately. \gls{x-bap} is a promising solution to this problem by providing accurate photometric colors using \gls{psf}-corrected aperture photometry.
